\documentclass[aspectratio=169,11pt]{beamer}
\usepackage[utf8]{inputenc}
\usepackage[T1]{fontenc}
\usepackage[spanish,es-nodecimaldot]{babel}
\usepackage{lmodern}
\usepackage{amsmath,amssymb}
\usepackage{tikz}
\usepackage{xcolor}
\usepackage{siunitx}
\usetikzlibrary{arrows.meta,calc,positioning,angles,quotes}
\sisetup{output-decimal-marker={,},per-mode=symbol}
\definecolor{azul}{RGB}{29,64,109}
\definecolor{morado}{RGB}{111,45,158}
\definecolor{verde}{RGB}{24,126,76}
\definecolor{rojo}{RGB}{185,55,60}
\definecolor{naranja}{RGB}{216,122,27}
\usetheme{Madrid}
\setbeamercolor{structure}{fg=azul}
\setbeamercolor{frametitle}{bg=azul,fg=white}
\setbeamercolor{block title}{bg=azul,fg=white}
\setbeamercolor{block body}{bg=azul!5}
\setbeamercolor{alerted text}{fg=morado}
\setbeamertemplate{navigation symbols}{}
\setbeamertemplate{footline}[frame number]
\setlength{\parskip}{3pt}
\newcommand{\res}[1]{\begin{center}\setlength{\fboxsep}{8pt}\fcolorbox{morado}{morado!5}{\color{morado}\bfseries\ensuremath{#1}}\end{center}}
\newcommand{\dato}[1]{\textbf{Datos:} #1}
\tikzset{force/.style={-{Stealth[length=3mm]},thick},axis/.style={-{Stealth[length=2mm]},gray,dashed},block/.style={draw=azul,fill=azul!12,thick,rounded corners=1pt},ground/.style={line width=2pt,gray!65}}
\title{Dinámica: leyes de Newton}
\subtitle{Diagramas de cuerpo libre, fricción, planos inclinados y poleas}
\author{Julián Arias Meza}
\institute{Física -- Educación Diversificada}
\date{}
\begin{document}
\begin{frame}\titlepage\end{frame}
\begin{frame}{Objetivos de aprendizaje}
\begin{enumerate}
\item Reconocer las fuerzas que actúan sobre un cuerpo.
\item Construir diagramas de cuerpo libre (DCL).
\item Descomponer fuerzas en componentes horizontal y vertical.
\item Calcular peso, normal y fuerza de fricción.
\item Aplicar la segunda ley de Newton a bloques y planos inclinados.
\item Determinar aceleración y tensión en sistemas de poleas.
\end{enumerate}
\end{frame}
\section{Fundamentos}
\begin{frame}{¿Qué estudia la dinámica?}
La dinámica relaciona las fuerzas que actúan sobre un cuerpo con los cambios en su movimiento.
\begin{block}{Fuerza}
Interacción capaz de cambiar la velocidad de un cuerpo o deformarlo. Es una magnitud vectorial.
\end{block}
\[1\ \mathrm{N}=1\ \mathrm{kg\,m/s^2}\]
\begin{block}{Peso}
\[P=mg,\qquad g=9{,}8\ \mathrm{m/s^2}\]
El peso es una fuerza y se expresa en newtons; la masa se expresa en kilogramos.
\end{block}
\end{frame}
\begin{frame}{Las tres leyes de Newton}
\begin{block}{Primera ley: inercia}
Si la fuerza resultante es cero, la velocidad permanece constante (incluido el reposo).
\[\sum\vec F=\vec 0\]
\end{block}
\begin{block}{Segunda ley: dinámica}
\[\boxed{\sum\vec F=m\vec a}\]
\end{block}
\begin{block}{Tercera ley: acción y reacción}
Las fuerzas de interacción entre dos cuerpos tienen igual magnitud y sentidos opuestos, pero actúan sobre cuerpos diferentes.
\end{block}
\end{frame}
\begin{frame}{Segunda ley en dos ejes}
\[\boxed{\sum F_x=ma_x}\qquad\boxed{\sum F_y=ma_y}\]
\begin{itemize}
\item Elegimos el sentido positivo antes de sumar fuerzas.
\item Una fuerza dirigida en sentido negativo se resta.
\item Si no hay aceleración vertical: \(\sum F_y=0\).
\item Si la fuerza resultante es cero, la aceleración es cero.
\end{itemize}
\begin{alertblock}{Cuidado}
La aceleración no es una fuerza y no se dibuja como fuerza en un DCL.
\end{alertblock}
\end{frame}
\begin{frame}{Fuerzas habituales}
\begin{columns}[T]
\column{.53\textwidth}
\begin{itemize}
\item \textbf{Peso} \(P=mg\): vertical hacia abajo.
\item \textbf{Normal} \(N\): perpendicular a la superficie.
\item \textbf{Tensión} \(T\): a lo largo de la cuerda.
\item \textbf{Fuerza aplicada} \(F\): dirección indicada.
\item \textbf{Fricción} \(f\): se opone al deslizamiento relativo o a su tendencia.
\end{itemize}
\column{.47\textwidth}
\centering
\begin{tikzpicture}[scale=.82]
\draw[ground] (-2,-.65)--(2,-.65);
\draw[block] (-.75,-.65) rectangle (.75,.45);
\draw[force,verde] (0,.45)--(0,1.65) node[above] {$\vec N$};
\draw[force,rojo] (0,-.15)--(0,-1.8) node[below] {$\vec P$};
\draw[force,azul] (.75,0)--(2,0) node[right] {$\vec F$};
\draw[force,morado] (-.75,-.35)--(-2,-.35) node[left] {$\vec f$};
\end{tikzpicture}
\end{columns}
\end{frame}
\begin{frame}{Cómo elaborar un diagrama de cuerpo libre}
\begin{enumerate}
\item Aislar el cuerpo que se estudia.
\item Elegir ejes coordenados convenientes.
\item Dibujar únicamente las fuerzas externas sobre ese cuerpo.
\item Asignar sentidos positivos y negativos.
\item Descomponer las fuerzas inclinadas cuando corresponda.
\end{enumerate}
\begin{center}
\begin{tikzpicture}[scale=.7]
\fill (0,0) circle (2pt);
\draw[force,verde] (0,0)--(0,1.4) node[above] {$N$};
\draw[force,rojo] (0,0)--(0,-1.4) node[below] {$P$};
\draw[force,azul] (0,0)--(1.8,0) node[right] {$F$};
\draw[force,morado] (0,0)--(-1.8,0) node[left] {$f$};
\end{tikzpicture}
\end{center}
\end{frame}
\begin{frame}{Descomposición de una fuerza inclinada}
\begin{columns}
\column{.53\textwidth}
Si \(\theta\) se mide desde la horizontal:
\[\boxed{F_x=F\cos\theta}\]
\[\boxed{F_y=F\operatorname{sen}\theta}\]
Los signos se asignan según el sentido real de cada componente.
\column{.47\textwidth}
\centering
\begin{tikzpicture}[scale=.92]
\draw[axis] (-.2,0)--(3.2,0) node[right] {$x$};
\draw[axis] (0,-.2)--(0,2.5) node[above] {$y$};
\draw[force,verde] (0,0)--(2.5,1.75) node[above] {$F$};
\draw[force,azul] (0,0)--(2.5,0) node[below] {$F_x$};
\draw[force,naranja] (2.5,0)--(2.5,1.75) node[right] {$F_y$};
\draw (0.75,0) arc[start angle=0,end angle=35,radius=.75];
\node at (1.05,.3) {$\theta$};
\end{tikzpicture}
\end{columns}
\end{frame}
\section{Bloques horizontales}
\begin{frame}{Ejemplo 1. Bloque sin fricción: planteamiento}
Un bloque de \(10\ \mathrm{kg}\) recibe una fuerza horizontal de \(50\ \mathrm{N}\). Determine su aceleración.
\begin{center}
\begin{tikzpicture}[scale=.85]
\draw[ground] (-2,-.55)--(2.6,-.55);
\draw[block] (-.8,-.55) rectangle (.8,.65);
\node at (0,.05) {$10\ \mathrm{kg}$};
\draw[force,azul] (.8,.05)--(2.7,.05) node[right] {$50\ \mathrm{N}$};
\draw[force,verde] (0,.65)--(0,1.65) node[above] {$N$};
\draw[force,rojo] (0,-.1)--(0,-1.5) node[below] {$P$};
\end{tikzpicture}
\end{center}
\dato{\(m=10\ \mathrm{kg}\), \(F=50\ \mathrm{N}\), \(f=0\).}
\end{frame}
\begin{frame}{Ejemplo 1. Solución}
\textbf{Fórmula}
\[\sum F_x=ma\]
\textbf{Sustitución}
\[50\ \mathrm{N}=(10\ \mathrm{kg})a\]
\textbf{Procedimiento}
\[a=\frac{50\ \mathrm{N}}{10\ \mathrm{kg}}\]
\textbf{Respuesta}
\res{a=5\ \mathrm{m/s^2}}
El bloque acelera hacia la derecha.
\end{frame}
\begin{frame}{Fuerza normal en una superficie horizontal}
Si el bloque permanece apoyado y \(a_y=0\):
\begin{block}{Sin fuerzas verticales adicionales}
\[\boxed{N=mg}\]
\end{block}
\begin{block}{Una fuerza tira hacia arriba con ángulo \(\theta\)}
\[\boxed{N=mg-F\operatorname{sen}\theta}\]
\end{block}
\begin{block}{Una fuerza empuja hacia abajo con ángulo \(\theta\)}
\[\boxed{N=mg+F\operatorname{sen}\theta}\]
\end{block}
Si \(F\operatorname{sen}\theta\ge mg\), puede perderse el contacto y \(N=0\).
\end{frame}
\begin{frame}{Fricción estática y cinética}
\begin{block}{Fricción estática: antes de deslizar}
\[\boxed{0\le f_s\le \mu_sN}\]
La fuerza estática se ajusta hasta un máximo \(f_{s,\max}=\mu_sN\).
\end{block}
\begin{block}{Fricción cinética: mientras desliza}
\[\boxed{f_k=\mu_kN}\]
\end{block}
\begin{itemize}
\item Los coeficientes \(\mu_s\) y \(\mu_k\) no tienen unidades.
\item En los ejemplos con movimiento se utiliza \(f_k\).
\end{itemize}
\end{frame}
\begin{frame}{Ejemplo 2. Tres fuerzas y fricción: enunciado}
Un bloque de \(40\ \mathrm{kg}\) se mueve sobre una superficie horizontal. Actúan \(F_1=80\ \mathrm{N}\) a la derecha, \(F_2=130\ \mathrm{N}\) a \(30^\circ\) sobre la horizontal y \(F_3=10\ \mathrm{N}\) a la izquierda. \(\mu_k=0{,}16\).
\begin{center}
\begin{tikzpicture}[scale=.72]
\draw[ground] (-3,-.7)--(4,-.7);
\draw[block] (-1,-.7) rectangle (1,.55);
\node at (0,-.08) {$40\ \mathrm{kg}$};
\draw[force,azul] (1,-.4)--(3.4,-.4) node[right] {$F_1$};
\draw[force,verde] (1,.35)--(3,1.5) node[right] {$F_2$};
\draw[force,rojo] (-1,0)--(-2.7,0) node[left] {$F_3$};
\draw[force,morado] (-1,-.5)--(-2.5,-.5) node[left] {$f_k$};
\draw (1.75,.35) arc[start angle=0,end angle=30,radius=.75];
\node at (2.2,.6) {$30^\circ$};
\end{tikzpicture}
\end{center}
Determine la normal, la fricción y la aceleración.
\end{frame}
\begin{frame}{Ejemplo 2. Normal}
\textbf{Datos}
\[m=40\ \mathrm{kg},\quad F_2=130\ \mathrm{N},\quad \theta=30^\circ\]
\textbf{Fórmulas}
\[P=mg,\qquad \sum F_y=0\]
\[N+F_2\operatorname{sen}\theta-P=0\]
\textbf{Sustitución y procedimiento}
\[P=40(9{,}8)=392\ \mathrm{N}\]
\[N=392-130\operatorname{sen}30^\circ=392-65\]
\res{N=327\ \mathrm{N}}
\end{frame}
\begin{frame}{Ejemplo 2. Fricción}
\textbf{Datos}
\[\mu_k=0{,}16,\qquad N=327\ \mathrm{N}\]
\textbf{Fórmula}
\[f_k=\mu_kN\]
\textbf{Sustitución}
\[f_k=0{,}16(327\ \mathrm{N})\]
\textbf{Procedimiento}
\[f_k=52{,}32\ \mathrm{N}\]
\res{f_k=52{,}32\ \mathrm{N}}
La fricción actúa hacia la izquierda, opuesta al movimiento.
\end{frame}
\begin{frame}{Ejemplo 2. Aceleración}
\textbf{Fórmula}
\[\sum F_x=ma\]
\textbf{Sustitución}
\[80+130\cos30^\circ-10-52{,}32=40a\]
\textbf{Procedimiento}
\[80+112{,}58-10-52{,}32=40a\]
\[130{,}26=40a\]
\[a=\frac{130{,}26}{40}\]
\res{a\approx3{,}26\ \mathrm{m/s^2}}
\end{frame}
\section{Planos inclinados}
\begin{frame}{Plano inclinado: ejes y componentes del peso}
Elegimos el eje \(x\) paralelo al plano y el eje \(y\) perpendicular al plano.
\begin{columns}
\column{.5\textwidth}
\[\boxed{P_x=mg\operatorname{sen}\theta}\]
\[\boxed{P_y=mg\cos\theta}\]
Sin otras fuerzas perpendiculares:
\[\boxed{N=mg\cos\theta}\]
\column{.5\textwidth}
\centering
\begin{tikzpicture}[scale=.85]
\begin{scope}[rotate=25]
\draw[ground] (-2.2,0)--(2.5,0);
\draw[block] (-.7,0) rectangle (.7,.85);
\coordinate (C) at (0,.43);
\draw[force,verde] (C)--++(0,1.35) node[above] {$N$};
\draw[force,azul] (C)--++(-1.5,0) node[left] {$P_x$};
\draw[force,morado] (C)--++(0,-1.15) node[right] {$P_y$};
\end{scope}
\draw[force,rojo] (-0.2,.39)--(0,-1.5) node[below] {$P$};
\end{tikzpicture}
\end{columns}
\end{frame}
\begin{frame}{Plano inclinado: casos frecuentes}
\begin{block}{Sin fricción ni fuerza aplicada (desciende)}
\[mg\operatorname{sen}\theta=ma\quad\Rightarrow\quad\boxed{a=g\operatorname{sen}\theta}\]
\end{block}
\begin{block}{Deslizamiento hacia abajo con fricción cinética}
\[mg\operatorname{sen}\theta-f_k=ma\]
\[\boxed{a=g(\operatorname{sen}\theta-\mu_k\cos\theta)}\]
\end{block}
\begin{alertblock}{Atención}
La fricción apunta hacia arriba del plano si el bloque desliza hacia abajo; cambia de sentido si el bloque desliza hacia arriba.
\end{alertblock}
\end{frame}
\begin{frame}{Ejemplo 3. Fuerza paralela al plano: enunciado}
Un bloque de \(12\ \mathrm{kg}\) está sobre un plano de \(25^\circ\), sin fricción. Se aplica \(F=100\ \mathrm{N}\) hacia arriba del plano.
\begin{center}
\begin{tikzpicture}[scale=.82]
\begin{scope}[rotate=25]
\draw[ground] (-2.6,0)--(3.2,0);
\draw[block] (-.75,0) rectangle (.75,.9);
\node at (0,.45) {$12\ \mathrm{kg}$};
\draw[force,verde] (.75,.45)--(2.8,.45) node[above] {$100\ \mathrm{N}$};
\end{scope}
\draw[dashed,gray] (-2.4,-1)--(2.8,-1);
\node at (-0.5,-.75) {$25^\circ$};
\end{tikzpicture}
\end{center}
Calcule \(P_x\), \(N\), fuerza resultante y aceleración.
\end{frame}
\begin{frame}{Ejemplo 3. Componente paralela del peso}
\textbf{Datos}
\[m=12\ \mathrm{kg},\quad \theta=25^\circ,\quad g=9{,}8\ \mathrm{m/s^2}\]
\textbf{Fórmula}
\[P_x=mg\operatorname{sen}\theta\]
\textbf{Sustitución}
\[P_x=12(9{,}8)\operatorname{sen}25^\circ\]
\textbf{Procedimiento}
\[P_x=117{,}6\operatorname{sen}25^\circ\]
\res{P_x\approx49{,}70\ \mathrm{N}}
Esta componente apunta hacia abajo del plano.
\end{frame}
\begin{frame}{Ejemplo 3. Fuerza normal}
\textbf{Fórmula}
\[\sum F_y=0\quad\Rightarrow\quad N-P_y=0\]
\[N=mg\cos\theta\]
\textbf{Sustitución}
\[N=12(9{,}8)\cos25^\circ\]
\textbf{Procedimiento}
\[N=117{,}6\cos25^\circ\]
\res{N\approx106{,}58\ \mathrm{N}}
\end{frame}
\begin{frame}{Ejemplo 3. Resultante y aceleración}
Positivo: hacia arriba del plano.
\textbf{Fórmula}
\[\sum F_x=ma\quad\Rightarrow\quad F-P_x=ma\]
\textbf{Sustitución}
\[100-49{,}70=12a\]
\textbf{Procedimiento}
\[F_R=50{,}30\ \mathrm{N},\qquad a=\frac{50{,}30}{12}\]
\res{F_R\approx50{,}30\ \mathrm{N}\qquad a\approx4{,}19\ \mathrm{m/s^2}}
El bloque acelera hacia arriba del plano.
\end{frame}
\section{Poleas}
\begin{frame}{Modelo de cuerda y polea ideales}
\begin{itemize}
\item La cuerda tiene masa despreciable y no se estira.
\item La polea es ideal: sin masa ni fricción.
\item La tensión tiene igual magnitud en los tramos de la misma cuerda.
\item Los cuerpos unidos comparten la magnitud de la aceleración.
\item Se escribe \(\sum F=ma\) para cada cuerpo por separado.
\end{itemize}
\begin{alertblock}{Regla esencial}
El peso de la masa colgante no es, en general, igual a la tensión cuando existe aceleración.
\end{alertblock}
\end{frame}
\begin{frame}{Máquina de Atwood: diagrama}
\begin{columns}
\column{.5\textwidth}
\centering
\begin{tikzpicture}[scale=.9]
\draw[thick] (-1.2,2.45)--(1.2,2.45);
\draw[thick] (0,2.45)--(0,2.13);
\draw[fill=gray!20] (0,1.78) circle (.35);
\draw[thick] (-.35,1.78)--(-.35,-.35);
\draw[thick] (.35,1.78)--(.35,.25);
\draw[block] (-.83,-1.05) rectangle (.13,-.35);
\draw[block,fill=green!12] (.08,-.45) rectangle (1.04,.25);
\node at (-.35,-.7) {$m_1$};\node at (.56,-.1) {$m_2$};
\draw[force,azul] (-1.15,-.2)--(-1.15,-1.2) node[below] {$a$};
\draw[force,verde] (1.35,-.4)--(1.35,.6) node[above] {$a$};
\end{tikzpicture}
\column{.5\textwidth}
Si \(m_1>m_2\), la masa mayor baja.
\[P_1-T=m_1a\]
\[T-P_2=m_2a\]
Al sumar:
\[P_1-P_2=(m_1+m_2)a\]
\end{columns}
\end{frame}
\begin{frame}{Máquina de Atwood: fórmulas de los apuntes}
\begin{block}{Aceleración}
\[\boxed{a=\frac{P_{\mathrm{mayor}}-P_{\mathrm{menor}}}{m_1+m_2}}\]
\end{block}
\begin{block}{Tensión usando la masa menor}
\[\boxed{T=P_{\mathrm{menor}}+m_{\mathrm{menor}}a}\]
\end{block}
\begin{block}{Tensión usando la masa mayor}
\[\boxed{T=P_{\mathrm{mayor}}-m_{\mathrm{mayor}}a}\]
\end{block}
Las dos expresiones de tensión deben dar el mismo resultado.
\end{frame}
\begin{frame}{Ejemplo 4. Máquina de Atwood: datos y aceleración}
Dos masas de \(12\ \mathrm{kg}\) y \(10\ \mathrm{kg}\) están conectadas por una polea ideal. Determine \(a\) y \(T\).
\textbf{Datos}
\[m_1=12\ \mathrm{kg},\quad m_2=10\ \mathrm{kg},\quad g=9{,}8\ \mathrm{m/s^2}\]
\textbf{Fórmula}
\[a=\frac{P_{\mathrm{mayor}}-P_{\mathrm{menor}}}{m_1+m_2}\]
\textbf{Sustitución}
\[a=\frac{12(9{,}8)-10(9{,}8)}{12+10}\]
\textbf{Procedimiento}
\[a=\frac{117{,}6-98}{22}=\frac{19{,}6}{22}\]
\res{a\approx0{,}8909\ \mathrm{m/s^2}}
\end{frame}
\begin{frame}{Ejemplo 4. Tensión por la masa menor}
\textbf{Fórmula}
\[T=P_{\mathrm{menor}}+m_{\mathrm{menor}}a\]
\textbf{Sustitución}
\[T=98+10\left(\frac{19{,}6}{22}\right)\]
\textbf{Procedimiento}
\[T=98+8{,}9091\]
\res{T\approx106{,}91\ \mathrm{N}}
La masa de \(12\ \mathrm{kg}\) desciende y la de \(10\ \mathrm{kg}\) asciende.
\end{frame}
\begin{frame}{Ejemplo 4. Comprobación con la masa mayor}
\textbf{Fórmula}
\[T=P_{\mathrm{mayor}}-m_{\mathrm{mayor}}a\]
\textbf{Sustitución}
\[T=117{,}6-12\left(\frac{19{,}6}{22}\right)\]
\textbf{Procedimiento}
\[T=117{,}6-10{,}6909\]
\res{T\approx106{,}91\ \mathrm{N}}
Se obtiene la misma tensión por ambos cuerpos.
\end{frame}
\begin{frame}{Bloque horizontal y masa colgante: esquema}
\begin{center}
\begin{tikzpicture}[scale=.9]
\draw[ground] (-3,0)--(1.6,0);
\draw[block] (-2.15,0) rectangle (-.9,.85);
\node at (-1.525,.42) {$m_1$};
\draw[fill=gray!20] (1.6,.47) circle (.34);
\draw[thick] (-.9,.67)--(1.6,.81);
\draw[thick] (1.94,.47)--(1.94,-1.1);
\draw[block,fill=green!12] (1.48,-1.9) rectangle (2.4,-1.1);
\node at (1.94,-1.5) {$m_2$};
\draw[force,azul] (-1.5,1.1)--(-.45,1.1) node[right] {$a$};
\draw[force,azul] (2.75,-1.1)--(2.75,-1.9) node[below] {$a$};
\draw[force,morado] (-2.15,.25)--(-3,.25) node[left] {$f_k$};
\end{tikzpicture}
\end{center}
La masa colgante desciende y el bloque horizontal se mueve hacia la polea. La fricción, si existe, se opone a ese desplazamiento.
\end{frame}
\begin{frame}{Sistema horizontal: diagramas de cuerpo libre}
\begin{columns}
\column{.5\textwidth}
\centering\textbf{Bloque horizontal \(m_1\)}
\begin{tikzpicture}[scale=.82]
\fill (0,0) circle (2pt);
\draw[force,verde] (0,0)--(0,1.5) node[above] {$N$};
\draw[force,rojo] (0,0)--(0,-1.5) node[below] {$P_1$};
\draw[force,azul] (0,0)--(1.7,0) node[right] {$T$};
\draw[force,morado] (0,0)--(-1.7,0) node[left] {$f_k$};
\end{tikzpicture}
\[T-f_k=m_1a\]
\column{.5\textwidth}
\centering\textbf{Masa colgante \(m_2\)}
\begin{tikzpicture}[scale=.82]
\fill (0,0) circle (2pt);
\draw[force,verde] (0,0)--(0,1.5) node[above] {$T$};
\draw[force,rojo] (0,0)--(0,-1.5) node[below] {$P_2$};
\end{tikzpicture}
\[P_2-T=m_2a\]
\end{columns}
\end{frame}
\begin{frame}{Sistema horizontal: fórmulas de los apuntes}
\begin{block}{Aceleración (masa colgante descendiendo)}
\[\boxed{a=\frac{P_{\mathrm{colgante}}-f_k}{m_1+m_2}}\]
\end{block}
\begin{block}{Tensión desde la masa del plano}
\[\boxed{T=m_{\mathrm{plano}}a+f_k}\]
\end{block}
\begin{block}{Tensión desde la masa colgante}
\[\boxed{T=P_{\mathrm{colgante}}-m_{\mathrm{colgante}}a}\]
\end{block}
Si no hay fricción, se utiliza \(f_k=0\).
\end{frame}
\begin{frame}{Ejemplo 5. Sistema horizontal sin fricción}
Un bloque de \(10\ \mathrm{kg}\) se conecta mediante una polea ideal a una masa colgante de \(15\ \mathrm{kg}\). Determine \(a\) y \(T\).
\textbf{Datos}
\[m_1=10\ \mathrm{kg},\quad m_2=15\ \mathrm{kg},\quad f_k=0\]
\textbf{Fórmula}
\[a=\frac{P_{\mathrm{colgante}}}{m_1+m_2}\]
\textbf{Sustitución}
\[a=\frac{15(9{,}8)}{10+15}\]
\textbf{Procedimiento}
\[a=\frac{147}{25}\]
\res{a=5{,}88\ \mathrm{m/s^2}}
\end{frame}
\begin{frame}{Ejemplo 5. Tensión}
\textbf{Fórmula}
\[T=m_{\mathrm{plano}}a\]
\textbf{Sustitución}
\[T=10(5{,}88)\]
\textbf{Procedimiento}
\[T=58{,}8\ \mathrm{N}\]
\res{T=58{,}8\ \mathrm{N}}
\textbf{Comprobación:}
\[T=147-15(5{,}88)=58{,}8\ \mathrm{N}\]
\end{frame}
\begin{frame}{Ejemplo 6. Hallar el coeficiente de fricción}
Un bloque de \(20\ \mathrm{kg}\) se conecta a una masa colgante de \(15\ \mathrm{kg}\). El sistema acelera a \(3\ \mathrm{m/s^2}\), con la masa colgante descendiendo.
\begin{block}{Determine}
\begin{enumerate}
\item La fuerza de fricción cinética.
\item El coeficiente de fricción cinética.
\item La tensión de la cuerda.
\end{enumerate}
\end{block}
\dato{\(m_1=20\ \mathrm{kg}\), \(m_2=15\ \mathrm{kg}\), \(a=3\ \mathrm{m/s^2}\).}
\end{frame}
\begin{frame}{Ejemplo 6. Fuerza de fricción}
\textbf{Fórmula}
\[a=\frac{P_{\mathrm{colgante}}-f_k}{m_1+m_2}\]
\textbf{Despeje}
\[f_k=P_{\mathrm{colgante}}-(m_1+m_2)a\]
\textbf{Sustitución}
\[f_k=15(9{,}8)-(20+15)(3)\]
\textbf{Procedimiento}
\[f_k=147-105\]
\res{f_k=42\ \mathrm{N}}
\end{frame}
\begin{frame}{Ejemplo 6. Coeficiente de fricción}
\textbf{Fórmulas}
\[N=m_1g,\qquad \mu_k=\frac{f_k}{N}\]
\textbf{Sustitución}
\[N=20(9{,}8)=196\ \mathrm{N}\]
\[\mu_k=\frac{42}{196}\]
\textbf{Procedimiento}
\[\mu_k\approx0{,}2143\]
\res{\mu_k\approx0{,}214}
El coeficiente de fricción no tiene unidades.
\end{frame}
\begin{frame}{Ejemplo 6. Tensión}
\textbf{Fórmula}
\[T=m_{\mathrm{plano}}a+f_k\]
\textbf{Sustitución}
\[T=20(3)+42\]
\textbf{Procedimiento}
\[T=60+42\]
\res{T=102\ \mathrm{N}}
\textbf{Comprobación:}
\[T=15(9{,}8)-15(3)=102\ \mathrm{N}\]
\end{frame}
\begin{frame}{Ejemplo 7. Ejercicio de los apuntes: enunciado}
Un bloque de \(20\ \mathrm{kg}\) se conecta a una masa colgante de \(33\ \mathrm{kg}\). La fuerza de fricción sobre el bloque horizontal es \(21{,}56\ \mathrm{N}\).
\begin{block}{Determine}
\begin{enumerate}
\item La aceleración del sistema.
\item La tensión de la cuerda.
\item Compruebe la tensión con ambos cuerpos.
\end{enumerate}
\end{block}
\dato{\(m_1=20\ \mathrm{kg}\), \(m_2=33\ \mathrm{kg}\), \(f_k=21{,}56\ \mathrm{N}\).}
\end{frame}
\begin{frame}{Ejemplo 7. Aceleración}
\textbf{Fórmula}
\[a=\frac{P_{\mathrm{colgante}}-f_k}{m_1+m_2}\]
\textbf{Sustitución}
\[a=\frac{33(9{,}8)-21{,}56}{20+33}\]
\textbf{Procedimiento}
\[a=\frac{323{,}4-21{,}56}{53}=\frac{301{,}84}{53}\]
\res{a\approx5{,}70\ \mathrm{m/s^2}}
\end{frame}
\begin{frame}{Ejemplo 7. Tensión desde el bloque horizontal}
\textbf{Fórmula}
\[T=m_{\mathrm{plano}}a+f_k\]
\textbf{Sustitución}
\[T=20\left(\frac{301{,}84}{53}\right)+21{,}56\]
\textbf{Procedimiento}
\[T\approx113{,}90+21{,}56\]
\res{T\approx135{,}46\ \mathrm{N}}
\end{frame}
\begin{frame}{Ejemplo 7. Comprobación desde la masa colgante}
\textbf{Fórmula}
\[T=P_{\mathrm{colgante}}-m_{\mathrm{colgante}}a\]
\textbf{Sustitución}
\[T=323{,}4-33\left(\frac{301{,}84}{53}\right)\]
\textbf{Procedimiento}
\[T\approx323{,}4-187{,}94\]
\res{T\approx135{,}46\ \mathrm{N}}
Se obtiene el mismo resultado sin redondear la aceleración antes del cálculo.
\end{frame}
\section{Resumen}
\begin{frame}{Formulario: bloques y planos}
\begin{columns}[T]
\column{.5\textwidth}
\textbf{Segunda ley}
\[\sum F_x=ma_x\]
\[\sum F_y=ma_y\]
\textbf{Peso}
\[P=mg\]
\textbf{Fuerza inclinada}
\[F_x=F\cos\theta\]
\[F_y=F\operatorname{sen}\theta\]
\column{.5\textwidth}
\textbf{Fricción cinética}
\[f_k=\mu_kN\]
\textbf{Plano inclinado}
\[P_x=mg\operatorname{sen}\theta\]
\[P_y=mg\cos\theta\]
\[N=mg\cos\theta\]
\end{columns}
\end{frame}
\begin{frame}{Formulario: poleas}
\textbf{Máquina de Atwood}
\[a=\frac{P_{\mathrm{mayor}}-P_{\mathrm{menor}}}{m_1+m_2}\]
\[T=P_{\mathrm{menor}}+m_{\mathrm{menor}}a\]
\vspace{.2cm}
\textbf{Bloque horizontal con masa colgante}
\[a=\frac{P_{\mathrm{colgante}}-f_k}{m_1+m_2}\]
\[T=m_{\mathrm{plano}}a+f_k\]
\[T=P_{\mathrm{colgante}}-m_{\mathrm{colgante}}a\]
\end{frame}
\begin{frame}{Estrategia para resolver cualquier ejercicio}
\begin{enumerate}
\item \textbf{Datos:} escribir masas, fuerzas, ángulos y coeficientes.
\item \textbf{DCL:} identificar todas las fuerzas de cada cuerpo.
\item \textbf{Ejes:} escoger sentidos positivos coherentes.
\item \textbf{Fórmulas:} escribir \(\sum F_x=ma_x\) y \(\sum F_y=ma_y\).
\item \textbf{Sustitución:} colocar los valores con unidades.
\item \textbf{Procedimiento:} despejar y calcular sin redondear antes de tiempo.
\item \textbf{Respuesta:} expresar magnitud, unidad y dirección.
\end{enumerate}
\end{frame}
\begin{frame}{Actividad de cierre}
Un bloque de \(15\ \mathrm{kg}\) se desliza sobre una superficie horizontal. Se le aplica una fuerza de \(100\ \mathrm{N}\) a \(30^\circ\) sobre la horizontal. El coeficiente de fricción cinética es \(0{,}20\).
\begin{block}{Calcule}
\begin{enumerate}
\item El peso del bloque.
\item La componente vertical de la fuerza aplicada.
\item La fuerza normal.
\item La fuerza de fricción cinética.
\item La aceleración horizontal.
\end{enumerate}
\end{block}
Dibuje el DCL y presente el procedimiento completo.
\end{frame}

% ==========================================
% DIAPOSITIVA FINAL
% ==========================================
\begin{frame}

\centering
\vfill

{\Huge\color{azul}
Fin de la presentación\par}

\vspace{0.7cm}

{\Large
Dinámica y segunda ley de Newton\par}

\vspace{0.8cm}

{\LARGE\color{morado}
\[
\sum \vec{F}=m\vec{a}
\]
}

\vfill

\end{frame}

\end{document}
